% id: 38
% book: 2026 정의여고 1-2 중간고사 프린트 · 창의 변형 1세트 · 01 평면좌표·직선·원·도형의 이동
% section: p13
% figure: none
% answer: ④
% answer_source: 계산
% variant_level: 2 (원본 전사)
% 이 파일은 build.mjs 가 items.json 에서 만든다. 고칠 때는 items.json 을 고친다.
\smash{\raisebox{1.5mm}{\dmkichul{2026 정의여고 1-2 중간고사 프린트 · 창의 변형 1세트 38번 · 13쪽}}}
좌표평면 위의 두 점~$\pt{A}(2{,}\mkern3mu 1)$, $\pt{B}(8{,}\mkern3mu 2)$와 $x$축 위의 점~$\pt{P}$에 대하여 $\seg{AP}+\seg{BP}$의 값이 최소가 되도록 하는 점~$\pt{P}$를 $\pt{P_0}$이라 하자. 직선~$\pt{BP_0}$을 $x$축에 대하여 대칭이동한 직선이 점~$(a{,}\mkern3mu -3)$을 지날 때, $a$의 값은?
\begin{choices32}{$4$}{$6$}{$8$}{$10$}{$12$}\end{choices32}
